\documentclass[a4paper]{article}     
\usepackage{amsfonts}
\usepackage{amsmath}
\usepackage{amssymb}
\usepackage{graphicx}
\newcommand{\bgcos}{\operatorname{Bgcos}}


\begin{document}

\noindent \large Auteur: Adsayan Gunaseelan \\
\noindent \large Studierichting: Informatica\\
\noindent \large Oefeningenreeks 2B nr. 1.8\\

\medskip

\normalsize
1.8) $\tan(\bgcos x)$, druk de volgende uitdrukking uit in $x$.\\

Stel $\bgcos x = y$ met $y \in [0,\pi]$\\

Dan, $\cos y = x$\\

$\tan y = \dfrac{\sin y}{\cos y}$\\

$=\dfrac{\sqrt{1-\cos^2 y}}{\cos y}$\\

$= \dfrac{\sqrt{1-x^2}}{x}$\\

$\sin^2 y + \cos^2 y = 1$\\

$\Leftrightarrow \sin y = \pm \sqrt{1-\cos^2 y}$\\

Omdat $y \in [0,\pi]$ nemen we $+$.\\

$\tan(\bgcos x) = \dfrac{\sqrt{1-x^2}}{x}$\\


\begin{thebibliography}{999}
%\bibitem{a} Referentie
\end{thebibliography}
\end{document} 